% Ad Atticum 13.46 --- headnote
% ad-atticum-13-46, 12 August 45 BC, Tusculum

\textit{Cicero to Atticus, written at the Tusculan villa on 12
August 45 BC --- Perseus dateline \textit{Scr. ia Tusculano prid.
Id. Sext. a. 709 (45)} (the \textit{ia} is a manuscript flicker
for \textit{in}). Four sections, the day's accumulated business
in one dispatch. It opens with one of Cicero's small Latin puns:
his freedman Pollex (``Thumb'') had promised to be there on the
Ides, and showed up at Lanuvium on the day before --- but proved
\textit{plane pollex, non index}, ``plainly a thumb, not a
pointer.'' Useless, in other words, for the errand he had been
sent on. The rest of the letter is Cicero piecing the day
together: an interview with Balbus at his Lanuvine villa,
brokered by Lepta; the freshly arrived letter from Caesar
promising to be back at Rome before the \textit{Ludi Romani};
and --- preserved here in a single sentence --- Caesar's report
of how he had read Cicero's \textit{Cato} and reread it until he
felt himself \textit{copiosior}, and how, by contrast, Brutus's
rival \textit{Cato} had made him think \textit{himself} eloquent.
A famous compliment, parried.}

\textit{The middle and end of the letter are estate business: the
formal sixty-day window for entering on the inheritance of the
deceased Cluvius is running, Vestorius (the Puteolan agent) is
being unhelpfully slow, and Pollex will have to be sent back to
make the entry on Cicero's behalf. Balbus is willing to write to
Caesar about the Cluvian gardens (\textit{nil liberalius}, ``nothing
could have been more obliging''), and supplies the figures from
Cluvius's will, including a legacy of fifty thousand sesterces to
Terentia. The closing line --- ``when this letter had already been
sealed, our courier arrived by night'' --- is the daily-letter
machinery showing through: Cicero seals, then breaks the seal to
add the late update, and Vestorius is at least partially absolved.
Register is hurried, businesslike, with the small Pollex pun and
the brief lament for Cossinius (\textit{dilexi hominem}, ``I loved
the man'') the only ornaments.}
